\chapter{Introduction générale}
\section{Contexte}
Le développement des outils numériques permet aux vendeurs et aux personnes intéressées par l’acquisition d’animaux de publier et de consulter des annonces à distance. En Tunisie, les besoins recouvrent à la fois les animaux d’élevage (ovins, caprins, bovins, équins, volailles) et les animaux de compagnie, dans le respect des règles de protection animale et des conditions légales de cession.

La variété des espèces soulève des questions de conception : tous les animaux ne disposent pas des mêmes caractéristiques pertinentes et toutes les annonces n’ont pas le même objectif. Une vente comporte un prix ; une adoption gratuite ou un don ne doivent pas dissimuler de contrepartie financière. Les alertes concernant un animal disparu répondent encore à un autre besoin et ne doivent pas être assimilées à des transactions commerciales.

\section{Problématique}
\begin{notebox}
\textbf{Comment concevoir une plateforme web tunisienne qui unifie les annonces de vente, d’adoption gratuite et de don pour plusieurs catégories d’animaux, propose une recherche multicritère fiable et offre un espace sécurisé pour les animaux perdus, présumés volés ou trouvés ?}
\end{notebox}

\section{Objectifs}
\begin{itemize}
 \item Définir un \textbf{référentiel cohérent} d’espèces, de races et de localisations.
 \item Concevoir des \textbf{annonces typées} avec règles de saisie adaptées au contexte métier.
 \item Permettre une \textbf{recherche combinée} par type d’offre, caractéristiques et localisation.
 \item Séparer \textbf{les annonces commerciales} des \textbf{alertes d’entraide}.
 \item Assurer confidentialité, \textbf{contrôle d’accès} et modération des informations sensibles.
 \item Réaliser progressivement le projet afin de maîtriser concrètement Spring Boot, Angular, les bases de données et les diagrammes UML.
\end{itemize}

\section{Démarche de travail}
Le projet suit une approche incrémentale. Chaque sprint associe \textbf{besoin métier}, \textbf{conception}, \textbf{développement personnel}, \textbf{tests} et \textbf{recette}. Le développeur écrit lui-même son code ; l’accompagnement consiste à expliquer les choix, analyser les erreurs et vérifier les critères d’acceptation. Le démarrage reste volontairement simple : le monolithe modulaire précède les technologies avancées et l’éventuelle extraction de microservices.

\section{Organisation du rapport}
Le chapitre~2 présente l’existant ; le chapitre~3 décrit la solution proposée et les règles métier ; le chapitre~4 recense le backlog ; le chapitre~5 organise les sprints ; les chapitres~6 et~7 détaillent des exemples de dossiers de conception avec diagrammes UML et maquettes ; les perspectives de réalisation sont ensuite précisées.
